% ================================================================================
% extra \newcommand's specific to this paper.
% ================================================================================

%% Oh no, we don't actually HAVE lambda!
%% Later we'll have to make this a higher-order calculus!
%% \newcommand{\ourcalc}[0]{\ensuremath{\lambda^{loc}}}
%% \newcommand{\lamadt}{\ensuremath{\lambda^{adt}}}
%% \newcommand{\lamcur}{\ensuremath{\lambda^{cur}}}

%% Proposal 1 - hey, we have capital lambda, and it's even the thing that binds
%% our lovely locations.
%% \newcommand{\ourcalc}[0]{\ensuremath{\Lambda^{loc}}}
%% \newcommand{\lamadt}{\ensuremath{\Lambda^{adt}}}
%% \newcommand{\lamcur}{\ensuremath{\Lambda^{cur}}}

% TODO: coming up with a new name.  RECONCILE usages with \ourcalc below.
% Ones using this macro have already been audited/reconciled.
%\newcommand{\sysname}[0]{}

%% %% Proposal 2 - Core, like GHC Core.
% \newcommand{\ourcalc}[0]{\ensuremath{Core^{loc}}}

%% Helpers
\newcommand{\pargibbon}[0]{Parallel Gibbon}

\newcommand{\ourcalc}{\rparcalc{}}
\newcommand{\rparcalc}{LoCal$^{\text{par}}$}
\newcommand{\parcalc}{LoCal$^{\text{par}}$}
\newcommand{\seqcalc}{LoCal}
% \newcommand{\lamadt} [0]{\ensuremath{Core^{adt}}}
\newcommand{\lamadt} [0]{HiCal\xspace}
% \newcommand{\lamcur} [0]{\ensuremath{Core^{cur}}}
% \newcommand{\lamcur} [0]{\textsf{Cursored}}
% \newcommand{\lamcur} [0]{\textsf{CursedCal}}

\newcommand{\lamcur}[0]{NoCal\xspace}

\newcommand{\gramdef}{\; ::= \;}
\newcommand{\gramor}{\; | \;}
\newcommand{\keywd}[1]{\mathit{#1}}
\newcommand{\gramwd}[1]{\texttt{#1}}
\newcommand{\sgramwd}[1]{\texttt{#1}}
\newcommand{\skeywd}[1]{\mathit{#1}}

%% Grammar
\newcommand{\PROG}{\keywd{top}}
\newcommand{\DD}{\keywd{dd}}
\newcommand{\VD}{\keywd{vd}}
\newcommand{\FD}{\keywd{fd}}
\newcommand{\DC}{\keywd{K}}
\newcommand{\sDC}{\skeywd{K}}
\newcommand{\TC}{\keywd{T}}
\newcommand{\EXPR}{\keywd{e}}
\newcommand{\sEXPR}{\skeywd{e}}
\newcommand{\DATA}{\gramwd{data}}
\newcommand{\TYP}{\keywd{\tau}}
\newcommand{\TYPC}{\keywd{\tau_c}}
\newcommand{\TYPIND}[1]{\gramwd{ind} \;  #1}
\newcommand{\hTYP}{\keywd{\hat{\tau}}}
\newcommand{\sTYP}{\skeywd{\tau}}
\newcommand{\var}{\svar}
\newcommand{\svar}{x}
\newcommand{\fvar}{\sfvar}
\newcommand{\sfvar}{f}
\newcommand{\yvar}{y}
\newcommand{\num}{n}
\newcommand{\ARROW}{\rightarrow}
\newcommand{\RP}{\keywd{dl}}
\newcommand{\sRP}{\skeywd{dl}}
\newcommand{\loc}{\skeywd{l}}
\newcommand{\sloc}{\skeywd{l}}
\newcommand{\concreteloc}[3]{\ensuremath{\langle #1, #2 \rangle ^{#3}}}
\newcommand{\reg}{\skeywd{r}}
\newcommand{\sreg}{\skeywd{r}}
\newcommand{\LS}{\keywd{ls}}
\newcommand{\LC}{\keywd{lc}}
\newcommand{\LE}{\keywd{le}}
\newcommand{\letpack}[3]{\gramwd{let}\;#1=#2\;\gramwd{in}\;#3}
\newcommand{\letloc}[3]{\gramwd{letloc}\;#1 = #2\;\gramwd{in}\;#3}
\newcommand{\letreg}[2]{\gramwd{letregion}\;#1\;\gramwd{in}\;#2}
\newcommand{\fapp}[2]{\fvar \;[#1]\; #2}
\newcommand{\TS}{\keywd{ts}}
\newcommand{\ssetbool}[4]{#4 = #3[#2\mapsto \keywd{#1}=\keywd{True}]}
\newcommand{\sinlinesetbool}[3]{#3[#2\mapsto \keywd{#1}=\keywd{True}]}
\newcommand{\CS}{\keywd{cs}}
\newcommand{\pat}{\keywd{pat}}
\newcommand{\ptr}[1]{\gramwd{ptr}\;\keywd{#1}}
\newcommand{\sind}{i}
\newcommand{\ind}{\keywd{i}}
\newcommand{\indj}{\keywd{j}}
\newcommand{\concreteind}[1]{#1}
\newcommand{\svind}[1]{#1\circ}
\newcommand{\svindplusoff}[2]{#1 + #2}
\newcommand{\indbef}[1]{#1\diamond}
\newcommand{\concretelocvar}{\keywd{cl}}
\newcommand{\fut}[1]{\text{fut}\;#1}
\newcommand{\concretelocbefore}[1]{\gramwd{before}\; #1}
\newcommand{\indivar}[1]{\gramwd{ivar}\; #1}
\newcommand{\ivarid}{\var}
\newcommand{\concretelocvardiamond}{\keywd{\concretelocvar\diamond}}
\newcommand{\indirection}[2]{\ensuremath{\text{\&}(#1,#2)}}
\newcommand{\indirectionvar}{hr}
\newcommand{\VAL}{\keywd{v}}
\newcommand{\ec}{\keywd{\mathcal{E}}}
\newcommand{\evalc}[1]{\ec \llbracket #1 \rrbracket }
\newcommand{\STOR}{\keywd{S}}
\newcommand{\stepsto}{\Rightarrow}
\newcommand{\parstepsto}{\Longrightarrow}
\newcommand{\rparstepsto}{\Longrightarrow_{rp}}
\newcommand{\costdyn}{\Downarrow}
\newcommand{\Name}{\Red{Name }}
\newcommand{\q}[1]{\texttt{#1}}
\newcommand{\frl}[1]{{\ensuremath \mathit{frl}(#1)}}
\newcommand{\caseclause}[2]{#1 \;\rightarrow \;#2}
\newcommand{\case}[2]{\gramwd{case}\; #1 \;\gramwd{of}\;#2}
\newcommand{\spat}{\keywd{spat}}
\newcommand{\switch}[2]{\gramwd{switch}\; #1 \;\gramwd{of}\;#2}
\newcommand{\readInt}[1]{\gramwd{\text{readInt}} \; #1}
\newcommand{\writeInt}[2]{\gramwd{\text{writeInt}} \; #1 \; #2}
\newcommand{\readTag}[1]{\gramwd{\text{readTag}} \; #1}
\newcommand{\writeTag}[2]{\gramwd{\text{writeTag}} \; #1 \; #2}
\newcommand{\readCursor}[1]{\gramwd{\text{readCursor}} \; #1}
\newcommand{\writeCursor}[2]{\gramwd{\text{writeCursor}} \; #1 \; #2}
\newcommand{\datacon}[3]{\ensuremath{#1 \;#2 \;#3}}
\newcommand{\litcon}[2]{\datacon{\keywd{L}}{#1}{#2}}
\newcommand{\litnum}{\keywd{n}}
\newcommand{\CP}{\keywd{cp}}
\newcommand{\tptr}[3]{\langle \mathit{ptr}\;#2\;#3 \rangle_{#1}}
\newcommand{\LM}{\keywd{M}}
\newcommand{\locis}[2]{\ensuremath{#1:#2}}
\newcommand{\lvar}{\keywd{lx}}
\newcommand{\has}[3]{\ensuremath{#1(#2) = #3}}
\newcommand{\startr}[1]{\gramwd{start}\; #1}
\newcommand{\afterl}[1]{\gramwd{after}\; #1}
\newcommand{\tightoverset}[2]{%
  \mathop{#2}\limits^{\vbox to -.5ex{\kern-0.75ex\hbox{$#1$}\vss}}}
\newcommand\set[1]{\{ \, \ensuremath{#1} \,\}}
\newcommand{\ecdot}{\bullet}
\newcommand{\heap}{\keywd{h}}
\newcommand{\heapval}{\keywd{hv}}
\newcommand{\alphaequiv}{=_{\alpha}}
\newcommand{\partup}{\ensuremath{\parallel}}

\newcommand{\tsetequivjudg}[2]{\ensuremath{#1 \approx_{ts} #2}}
\newcommand{\taskequivjudg}[3]{\ensuremath{#1; #2 \approx_{t} #3}}
\newcommand{\storequivjudg}[5]{\ensuremath{#1; (#2, #3) \approx_{s} (#4, #5)}}

% Cost graphs
\newcommand{\costgraphnode}{\loc}
\newcommand{\compgraphnodevar}{n}
\newcommand{\compgraphvar}{g}
\newcommand{\compgraphedgeset}{E}
\newcommand{\compgraphseqcompop}{\oplus}
\newcommand{\compgraphseqcomp}[2]{#1 \compgraphseqcompop #2}
\newcommand{\compgraphparcomp}[2]{#1 \otimes #2}
\newcommand{\heapgraph}{h}
\newcommand{\compgraphnode}[1]{[#1]}
\newcommand{\compgraph}[3]{(#1;#2;#3)}
\newcommand{\compgraphnodefresh}[1]{(#1 \; \mathbf{fresh})}


%% Environments
\newcommand{\LENV}{\keywd{L}}
\newcommand{\RENV}{\keywd{R}}
\newcommand{\CENV}{\keywd{C}}
\newcommand{\TENV}{\keywd{\Gamma}}
\newcommand{\EENV}{\keywd{E}}
\newcommand{\SENV}{\keywd{\Sigma}}
\newcommand{\MENV}{\keywd{M}}
\newcommand{\FENV}{\keywd{F}}
\newcommand{\AENV}{\keywd{A}}
\newcommand{\NENV}{\keywd{N}}

\newcommand{\lTENV}{\overharpoon{\keywd{\Gamma}}}
\newcommand{\lSENV}{\overharpoon{\keywd{\Sigma}}}
\newcommand{\lCENV}{\overharpoon{\keywd{C}}}
\newcommand{\lAENV}{\overharpoon{\keywd{A}}}
\newcommand{\lNENV}{\overharpoon{\keywd{N}}}

\newcommand{\liTENV}[1]{\overharpoon{\keywd{\Gamma}_{#1}}}
\newcommand{\liSENV}[1]{\overharpoon{\keywd{\Sigma}_{#1}}}
\newcommand{\liCENV}[1]{\overharpoon{\keywd{C}_{#1}}}
\newcommand{\liAENV}[1]{\overharpoon{\keywd{A}_{#1}}}
\newcommand{\liNENV}[1]{\overharpoon{\keywd{N}_{#1}}}

\newcommand{\TASKVAR}{\keywd{T}}
\newcommand{\SEQSTATE}{\keywd{t}}

\newcommand{\TASKSET}{\mathbb{T}}
\newcommand{\CENVMAP}{\mathbb{C}}
\newcommand{\AENVMAP}{\mathbb{A}}
\newcommand{\NENVMAP}{\mathbb{N}}
\newcommand{\SENVMAP}{\mathbbl{\Sigma}}
\newcommand{\TENVMAP}{\mathbbl{\Gamma}}

\newcommand{\singlewriter}[3]{\ensuremath{HasSingleWriter(#1, #2, #3)}}
\newcommand{\getwriter}[3]{\ensuremath{GetSingleWriter(#1, #2, #3)}}
\newcommand{\isval}[1]{\ensuremath{IsVal(#1)}}
\newcommand{\taskcomplete}[1]{\ensuremath{TaskComplete(#1)}}
\newcommand{\xor}{\ensuremath{\oplus}}

%% Loc helpers
\newcommand{\inloc}[0]{\ensuremath{\downarrow\sloc}}
\newcommand{\outloc}[0]{\ensuremath{\uparrow\sloc}}
\newcommand{\locreg}[2]{\ensuremath{{#1}^{#2}}}
\newcommand{\tyatlocreg}[3]{#1 \ensuremath{@} \locreg{#2}{#3}}
\newcommand{\styatlocreg}[3]{#1\ensuremath{@}\locreg{#2}{#3}}
\newcommand{\Cursorize}[0]{Cursorize}
\newcommand{\fresh}[0]{\keywd{fresh}}

%% Meta functions

% Takes FROM,TO, i.e. \subst{e}{x}{v}.
\newcommand{\subst}[3]{\ensuremath{#1[#3/#2]}}
\newcommand{\substtwo}[3]{\ensuremath{#1\llbracket#3/#2\rrbracket}}

\newcommand{\typeofcon}{\keywd{TypeOfCon}}
\newcommand{\typeoffield}{\keywd{TypeOfField}}
\newcommand{\kargtys}{\keywd{ArgTysOfConstructor}}
\newcommand{\allocptr}[2]{\keywd{MaxIdx}(#1,#2)}

% Update(M,key,val), e.g. M[key>val]
%\newcommand{\update}[3]{\ensuremath{((#2 \mapsto #3)\; #1)}}
\newcommand{\update}[3]{\ensuremath{#1+\{#2 \mapsto #3\}}}

\newcommand{\nfofi}{\keywd{Nf}}
\newcommand{\incri}{\keywd{Incr}}
\newcommand{\followinm}{\keywd{Deref}}
\newcommand{\rorngi}[2]{[#1, #2)}
\newcommand{\rorngh}{\keywd{RORngH}}
\newcommand{\storelookupone}[3]{\hat{#1}(#2,#3)}
\newcommand{\maplookupone}[2]{\hat{#1}(#2)}
\newcommand{\advancei}{\keywd{After}}
\newcommand{\mergeh}{\keywd{MergeH}}
\newcommand{\mergem}{\keywd{MergeM}}
\newcommand{\merges}{\keywd{MergeS}}
\newcommand{\tie}{\keywd{LinkFields}}
\newcommand{\linkfields}{\keywd{LinkFields}}

\newcommand{\ewitnesstext}{\ensuremath{\vdash_{\textit{ew}}}}
\newcommand{\ewitness}[4]{\ensuremath{#1;#2;#3 \ewitnesstext #4}}
\newcommand{\tietext}{\ensuremath{\vdash_{\textit{tie}}}}
\newcommand{\storewf}[7]{\ensuremath{#1;#2;#3;#4;#5 \vdash_{wf} #6;#7}}
\newcommand{\storewfcfa}[4]{\ensuremath{#1 \vdash_{wf_{cfc}} #3;#4}}
\newcommand{\storewfca}[5]{\ensuremath{#1;#2 \vdash_{wf_{ca}} #4;#5}}
\newcommand{\storewftasks}[1]{\ensuremath{\vdash_{wf_{tasks}} #1}}

\newcommand{\tcfun}{\ensuremath{\vdash_{fun}}}
\newcommand{\tcpat}{\ensuremath{\vdash_{pat}}}
\newcommand{\tcts}{\ensuremath{\vdash_{ts}}}
\newcommand{\tctask}{\ensuremath{\vdash_{task}}}
\newcommand{\tctaskset}{\ensuremath{\vdash_{taskset}}}

\newcommand{\parmergemenv}[2]{\mergem(#1, #2)}
\newcommand{\parmergestor}[2]{\merges(#1, #2)}
\newcommand{\parmergetask}[2]{\keywd{MergeT}(#1, #2)}

%% Special references for proofs.
\newcommand{\refwellformed}[2]{WF~\ref{#1};\ref{#2}}
\newcommand{\refendwitness}[2]{EW~\ref{#1};\ref{#2}}
\newcommand{\refts}[1]{T-#1}
\newcommand{\refcase}[1]{Case-\ref{#1}}
\newcommand{\elimexists}[0]{$\exists$ elim}
\newcommand{\refinst}[0]{Inst.}

\newcommand{\tfunctiondef}{T-Function-Definition}
\newcommand{\tdatacon}{T-DataConstructor}
\newcommand{\tdataconivars}{T-DataConstructor-Ivars}
\newcommand{\ddatacon}{D-DataConstructor}
\newcommand{\dletloctag}{D-LetLoc-Tag}
\newcommand{\dletlocafter}{D-LetLoc-After}
\newcommand{\dletlocafterpar}{D-Par-LetLoc-After}
\newcommand{\dletlocstart}{D-LetLoc-Start}
\newcommand{\dapp}{D-App}
\newcommand{\dletregion}{D-LetRegion}
\newcommand{\tcase}{T-Case}
\newcommand{\dcase}{D-Case}
\newcommand{\tpat}{T-Pattern}
\newcommand{\tprogram}{T-Program}
\newcommand{\tllafter}{T-LetLoc-After}
\newcommand{\tlltag}{T-LetLoc-Tag}
\newcommand{\tllstart}{T-LetLoc-Start}
\newcommand{\tlregion}{T-LetRegion}
\newcommand{\tvar}{T-Var}
\newcommand{\ttaskset}{T-Taskset}
\newcommand{\ttasksetempty}{T-Taskset-Empty}
\newcommand{\ttask}{T-Task}
\newcommand{\tlet}{T-Let}
\newcommand{\dletexp}{D-Let-Expr}
\newcommand{\dletval}{D-Let-Val}
\newcommand{\dletexppar}{D-Par-Let}
\newcommand{\dparcasejoin}{D-Par-Case-Join}
\newcommand{\tapp}{T-App}
\newcommand{\tconcreteloc}{T-Concrete-Loc}
\newcommand{\dparsteptask}{D-Par-Step}
\newcommand{\emptytenv}{\emptyset}

%% Region-Parallel transitions
\newcommand{\transformsto}{\longrightarrow}
\newcommand{\Indr}{I}
\newcommand{\trdatadecl}{Tr-DataDecl}
\newcommand{\dregparsteptask}{D-Par-Step}
\newcommand{\dregparletexppar}{D-Par-Let-Fork}
\newcommand{\dregparletexpseq}{D-Par-Let}
\newcommand{\dregparletlocafter}{D-Par-LetLoc-After}
\newcommand{\dregparletlocafternewreg}{D-LetLoc-After-NewReg}
\newcommand{\dregparcasejoin}{D-Par-Case-Join}
\newcommand{\dregparcase}{D-Par-Case}
\newcommand{\dregpardatacon}{D-Par-DataConstructor}
\newcommand{\dregpardataconjoin}{D-Par-DataConstructor-Join}
\newcommand{\dregparletlocstart}{D-Par-LetLoc-Start}
\newcommand{\dregparletloctag}{D-Par-LetLoc-Tag}
\newcommand{\dregparletexp}{D-Par-Let-Expr}
\newcommand{\dregparletval}{D-Par-Let-Val}
\newcommand{\dregparapp}{D-Par-App}
\newcommand{\dregparletregion}{D-Par-LetRegion}

%%
\newcommand{\wftask}{\ensuremath{\vdash_{wf_{task}}}}

%% Formatting

%% fdtools, harpoon are packages that provide this feature, but
%% have their own set of problems. We do this by hand instead:
%% https://tex.stackexchange.com/questions/304622
%% \makeatletter
%% \newcommand*\MY@rightharpoonupfill@{%
%%   \arrowfill@\relbar\relbar\rightharpoonup
%% }
%% \newcommand*\overrightharpoon{%
%%   \mathpalette{\overarrow@\MY@rightharpoonupfill@}%
%% }
%% \makeatother
\newcommand{\overharpoon}[1]{\overrightharpoon{#1}}

\newcommand{\taskconfig}[3]{(#1, #2, #3)}

%% \newcommand{\MPL}{\text{MPL}}
\newcommand{\MPL}{\text{MaPLe}}

\newcommand{\vsmaple}[1]{$\inferrule{\MPL{}}{\text{Ours}}{#1}$}
\newcommand{\vsocaml}[1]{$\inferrule{\text{OCaml}}{\text{Ours}}{#1}$}
\newcommand{\vsghc}[1]{$\inferrule{\text{GHC}}{\text{Ours}}{#1}$}
\newcommand{\buildfib}{buildtreeHvyLf}

\newcommand{\appendixref}[1]{appendix~\ref{#1}}
\newcommand{\Appendixref}[1]{Appendix~\ref{#1}}


\newcommand{\metafmerging}{
\[
\begin{array}{lcl}
  %% \merges(\STOR_1, \STOR_2) & = & \STOR_1 \cup \STOR_2 \\\\
  \merges(\STOR_1, \STOR_2) & = & \set{\reg \mapsto \mergeh(\heap_ 1, \heap_2) \mid (\reg \mapsto \heap_1) \in \STOR_1,  (\reg \mapsto \heap_2) \in \STOR_2} \\
                            & \cup & \set{\reg \mapsto \heap \mid (\reg \mapsto \heap) \in \STOR_1,  \reg \not\in dom(\STOR_2)} \\
                            & \cup & \set{\reg \mapsto \heap \mid (\reg \mapsto \heap) \in \STOR_2,  \reg \not\in dom(\STOR_1)} \\\\
  \mergeh(\heap_1, \heap_2) & = & \set{\concreteind{\ind_1} \mapsto \heapval \mid (\concreteind{\ind_1} \mapsto \heapval) \in \heap_1, (\concreteind{\ind_2} \mapsto \heapval) \in \heap_2, \concreteind{\ind_1} = \concreteind{\ind_2}} \\
                            & \cup & \set{\concreteind{\ind} \mapsto \heapval \mid (\concreteind{\ind} \mapsto \heapval) \in \heap_1, \concreteind{\ind} \not\in \set{\concreteind{\ind'} | \concreteind{\ind'} \in dom(\heap_2)}} \\
                            & \cup & \set{\concreteind{\ind} \mapsto \heapval \mid (\concreteind{\ind} \mapsto \heapval) \in \heap_2, \concreteind{\ind} \not\in \set{\concreteind{\ind'} | \concreteind{\ind'} \in dom(\heap_1)}} \\\\

  \mergem(\MENV_1, \MENV_2)
  %
  & = & \set{\locreg{\loc}{\reg} \mapsto \concreteloc{\reg}{\concreteind{\ind_1}}{} \mid
  (\locreg{\loc}{\reg} \mapsto \concreteloc{\reg}{\concreteind{\ind_1}}{}) \in \MENV_1,
  (\locreg{\loc}{\reg} \mapsto \concreteloc{\reg}{\concreteind{\ind_2}}{}) \in \MENV_2,
  \ind_1 = \ind_2} \\
  %
  & \cup & \set{\locreg{\loc}{\reg} \mapsto \concretelocvar \mid
  (\locreg{\loc}{\reg} \mapsto \concretelocvar) \in \MENV_1,
  \locreg{\loc}{\reg} \notin \MENV_2} \\
  %
  & \cup & \set{\locreg{\loc}{\reg} \mapsto \concretelocvar \mid
  \locreg{\loc}{\reg} \notin \MENV_1,
  (\locreg{\loc}{\reg} \mapsto \concretelocvar) \in \MENV_2} \\
  %
  & \cup & \set{\locreg{\loc}{\reg} \mapsto \concreteloc{\reg}{\indj}{} \mid
  (\locreg{\loc}{\reg} \mapsto \concreteloc{\reg}{\indivar{\ivarid}}{}) \in \MENV_1,
  (\locreg{\loc}{\reg} \mapsto \concreteloc{\reg}{\concreteind{\indj}}{}) \in \MENV_2} \\
  %
  & \cup & \set{\locreg{\loc}{\reg} \mapsto \concreteloc{\reg}{\indj}{} \mid
  (\locreg{\loc}{\reg} \mapsto \concreteloc{\reg}{\concreteind{\indj}}{}) \in \MENV_1,
  (\locreg{\loc}{\reg} \mapsto \concreteloc{\reg}{\indivar{\ivarid}}{}) \in \MENV_2} \\
  %
\end{array}
\]
}
